\section{Rigidity of the equality case}\label{sec:rigidity}
We prove Theorem~\ref{thm:equality}. Let $M^n$ and $E$ satisfy its
hypotheses, and assume equality in the Euclidean perimeter bound.
In particular, $E$ is bounded and has positive finite volume and
finite ambient perimeter. Retain the notation
\[
 m=n-1,\qquad V=\vol(E),\qquad
 R=(V/\omega_n)^{1/n},\qquad A=P(E)=s_mR^m.
\]
By \Cref{prop:equality-boundary}, the perimeter support is a compact set
$\Gamma=\Sigma\cup Z$, its regular part carries all the perimeter,
and its outward trace curvature is $H_0=m/R$.
We first show that $\Gamma$ is contained in the exponential image of a
Euclidean sphere of radius $R$. The argument depends on the dimension.
A single final argument then recovers both the set and the metric on its
interior.

\subsection{A barycenter and a flux inequality}
We record two observations used in the higher-dimensional and curve cases.
If $\mu$ is a nonzero finite positive measure of compact support in $M$,
the function
\[
 \mathcal F(p)=\int_M D_p(x)\,d\mu(x),\qquad
 D_p(x)=\tfrac12 d(p,x)^2,
\]
is continuous and proper. Indeed, if the support lies in $B_a(o)$, then
$d(p,x)\ge d(p,o)-a$, so $\mathcal F(p)\to\infty$ as $p$ leaves compact
sets. Completeness and Hopf--Rinow give a minimizer. Squared distance is
smooth on $M\times M$, and differentiation on compact sets gives
\begin{equation}\label{eq:barycenter}
 \int_M\log_p x\,d\mu(x)=0
\end{equation}
at any minimizer. Uniqueness of the minimizer is not needed.

Let $T$ be either $\Sigma$ or a compact boundaryless component of $\Sigma$,
and put $a=\area(T)$. The global first-variation identity from
\Cref{sec:equality-regularity} applies on $T$. The distance comparison
\Cref{lem:distance-comparison} yields, for every $p\in M$,
\begin{align}
 ma
 &\le \int_T \divg_{T\Sigma}\nabla D_p\,d\sigma
   =\frac mR\int_T\langle\nabla D_p,N\rangle\,d\sigma \notag\\
 &\le \frac mR\left(a\int_T r_p^2\,d\sigma\right)^{1/2}.
 \label{eq:flux-bound}
\end{align}
All fields may be multiplied by an ambient cutoff equal to one near
$\Gamma$, so their lack of compact support on $M$ causes no difficulty.
In particular,
\begin{equation}\label{eq:lower-moment}
 \int_T r_p^2\,d\sigma\ge R^2a.
\end{equation}
If the opposite inequality also holds, all inequalities in
\Cref{eq:flux-bound} are equalities. Consequently
\begin{equation}\label{eq:radial-normal}
 \nabla D_p=RN\quad\text{almost everywhere on }T.
\end{equation}
For example, the squared $L^2$ norm of the difference is
\[
 \int_T r_p^2\,d\sigma+R^2a
       -2R\int_T\langle\nabla D_p,N\rangle\,d\sigma=0.
\]
Continuity upgrades \Cref{eq:radial-normal} to every point of $T$:
a nonempty open regular patch has positive area. Since $|N|=1$,
this puts $T$ on the distance sphere $r_p=R$.

\subsection{Dimensions at least four}
Here $m\ge3$. The cutoff passage in \Cref{sec:equality-regularity}
has established \Cref{eq:equality-sobolev} for every smooth ambient
function restricted to $\Sigma$. Retain
$q=2m/(m-2)$ and $c=4/(m-2)=q-2$.
The constant function $v=1$ attains equality, since $A=s_mR^m$.

Let $\phi$ be such a restriction with
$\int_\Sigma\phi\,d\sigma=0$. Boundedness permits Taylor expansion for
$v=1+\tau\phi$:
\begin{equation}\label{eq:critical-expansion}
 \left(\int_\Sigma|1+\tau\phi|^q\,d\sigma\right)^{2/q}
 =A^{2/q}\left[1+(q-1)\frac{\tau^2}{A}
                \int_\Sigma\phi^2\,d\sigma+o(\tau^2)\right].
\end{equation}
The constant terms in \Cref{eq:equality-sobolev} agree, and the linear
terms vanish. Since $m s_m^{2/m}A^{2/q-1}=m/R^2$ and $c=q-2$,
comparison of the quadratic terms proves
\begin{equation}\label{eq:equality-spectral-gap}
 \int_\Sigma|\nabla_\Sigma\phi|^2\,d\sigma
 \ge \frac m{R^2}\int_\Sigma\phi^2\,d\sigma.
\end{equation}
The cutoff limit has already been taken before the expansion; no
mean-zero condition on cutoff approximations is asserted.

Choose a barycenter $p$ for the measure $d\sigma$ on $\Sigma$ and an
orthonormal basis $e_1,\ldots,e_n$ of $T_pM$.
The smooth functions
$\phi_i(x)=\langle\log_p x,e_i\rangle$ have zero mean by
\Cref{eq:barycenter}. Because $d\log_p$ is contracting,
\[
 \sum_{i=1}^n|\nabla_\Sigma\phi_i|^2
 =\sum_{j=1}^m|d\log_p(v_j)|^2\le m
\]
for any orthonormal basis $v_1,\ldots,v_m$ of $T_x\Sigma$.
Summing \Cref{eq:equality-spectral-gap} gives
\begin{equation}\label{eq:upper-moment}
 \int_\Sigma r_p^2\,d\sigma\le R^2A.
\end{equation}
Together with \Cref{eq:lower-moment}, this yields
\Cref{eq:radial-normal} on $\Sigma$. Its density in $\Gamma$ therefore
implies
\begin{equation}\label{eq:central-sphere}
 \Gamma\subset\{x\in M:d(p,x)=R\}.
\end{equation}
Neither connectedness nor completeness of the regular part was used.

\subsection{Dimension two}
Now $m=1$ and $\Gamma=\Sigma$ is a compact embedded $C^{1,1}$
curve without boundary, with total length $A=2\pi R$.
Choose a connected component $T$ and let $L=\operatorname{length}(T)>0$.
A connected compact one-dimensional manifold without boundary admits a
periodic arclength parametrization of period $L$. Here this can also be
seen by continuing an oriented arclength chart: a traverse of arbitrarily
large length cannot stay injective in a component of finite length, and
local uniqueness gives a least positive period. Its image is open and
closed in the component, and one least period traverses it once.

The periodic Wirtinger inequality says that
\begin{equation}\label{eq:wirtinger}
 \int_T |\nabla_T\phi|^2\,d\sigma
 \ge \left(\frac{2\pi}{L}\right)^2
       \int_T\phi^2\,d\sigma,
 \qquad \int_T\phi\,d\sigma=0.
\end{equation}
For completeness, in arclength coordinates the assertion follows by
expanding a trigonometric polynomial in Fourier modes; the zero mode is
absent. Periodic convolution approximates a $C^1$ function and its
derivative uniformly, preserving the zero mean, so the same inequality
holds for all the tests used below.

Take a barycenter of the area measure on $T$. Apply
\Cref{eq:wirtinger} to its logarithm coordinates and use the contraction
of $d\log_p$ as above. The result is
\[
 \int_T r_p^2\,d\sigma\le \frac{L^3}{(2\pi)^2}.
\]
On the other hand, \Cref{eq:lower-moment} gives
$\int_T r_p^2\,d\sigma\ge R^2L$. Thus $L\ge2\pi R=A$.
Since $L\le A$, this component has all the length, and every inequality
just used is an equality. Another component would contain a regular open
arc of positive length, which is impossible. Hence $T=\Gamma$, and
\Cref{eq:radial-normal} again proves \Cref{eq:central-sphere}.

\subsection{Dimension three}
Here $m=2$ and $\Sigma=\Gamma$ is a compact embedded $C^{1,1}$ surface,
with $A=4\pi R^2$. Write $h=1/R$, so its trace curvature is $2h$.
The critical Sobolev argument used for $m\ge3$ is replaced by a punctured
radial flux calculation.

Fix $p\in\Sigma$ and, away from $p$, define
\[
 r=r_p,\qquad Y=r^{-2}\nabla D_p=\frac{\nabla r}{r},
 \qquad t=\langle Y,N\rangle.
\]
Since $\Hess D_p\ge g$ and $|\nabla r|=1$,
\begin{align}
 \divg_{T\Sigma}Y
 &=r^{-2}\tr_{T\Sigma}\Hess D_p
      -2r^{-2}|\nabla_\Sigma r|^2 \notag\\
 &\ge2r^{-2}\bigl(1-|\nabla_\Sigma r|^2\bigr)=2t^2.
 \label{eq:puncture-divergence}
\end{align}
Choose a smooth nondecreasing function $\chi:[0,\infty)\to[0,1]$
that is zero on $[0,1]$ and one on $[2,\infty)$.
The field $\chi(r/\varepsilon)Y$ extends smoothly by zero near $p$.
Apply the global first-variation identity and differentiate the cutoff.
The term containing $\chi'$ is nonnegative, and
\Cref{eq:puncture-divergence} gives
\begin{equation}\label{eq:puncture-square}
 \int_\Sigma\chi(r/\varepsilon)
       \left[\frac{h^2}{2}-2\left(t-\frac h2\right)^2\right]d\sigma
 \ge J_\varepsilon,
 \quad
 J_\varepsilon=
 \int_\Sigma\frac{\chi'(r/\varepsilon)}{\varepsilon r}
                 |\nabla_\Sigma r|^2\,d\sigma.
\end{equation}
Indeed the left integrand before completing the square is
$2ht-2t^2=h^2/2-2(t-h/2)^2$.

We claim that
\begin{equation}\label{eq:annular-flux}
 \lim_{\varepsilon\downarrow0}J_\varepsilon=2\pi.
\end{equation}
Use exponential coordinates at $p$, with the first two coordinate axes
in $T_p\Sigma$. In a neighborhood of $p$ the entire support is a graph
$(y,b(y))$, where $b(0)=0$ and $Db(0)=0$.
For small $\varepsilon$ this one graph contains the whole annulus
contributing to $J_\varepsilon$; no other sheet or component enters it.
Set $y=\varepsilon z$. For fixed $z\ne0$,
\[
 \frac r\varepsilon\longrightarrow |z|,\qquad
 \frac{d\sigma}{\varepsilon^2}\longrightarrow dz,
 \qquad |\nabla_\Sigma r|\longrightarrow1.
\]
Here $r=|(y,b(y))|$ exactly, by the radial distance formula.
On the support of $\chi'$ one has $1\le r/\varepsilon\le2$.
Thus $z$ stays in a fixed bounded disk and, using the vanishing slope of
$b$, away from zero. The area density is uniformly bounded there,
$|\nabla_\Sigma r|\le1$, and
$|\chi'(r/\varepsilon)/(r/\varepsilon)|\le\|\chi'\|_\infty$.
Dominated convergence after extension by zero to the fixed disk gives
\[
 \lim_{\varepsilon\downarrow0}J_\varepsilon
 =\int_{\R^2}\frac{\chi'(|z|)}{|z|}\,dz
 =2\pi\int_1^2\chi'(s)\,ds=2\pi.
\]
This proves \Cref{eq:annular-flux}.

Because $h^2A/2=2\pi$, \Cref{eq:puncture-square} implies
\[
 0\le\int_\Sigma\chi(r/\varepsilon)(t-h/2)^2\,d\sigma
 \le \frac{h^2}{4}\int_\Sigma\chi(r/\varepsilon)\,d\sigma
              -\frac12J_\varepsilon\longrightarrow0.
\]
Fatou's lemma yields $t=h/2$ almost everywhere off $p$, without any
prior integrability assumption on $t^2$ at the pole. Continuity makes
this identity valid at every point of $\Sigma\setminus\{p\}$.

This everywhere off-diagonal identity has been proved for each fixed
$p$. We may therefore now fix $x_0\in\Sigma$ and evaluate it at $x_0$
for every $p\ne x_0$; there is no intersection of uncountably many
exceptional null sets. The identity
$\nabla D_p(x_0)=-\log_{x_0}p$ gives
\[
 -\langle\log_{x_0}p,N(x_0)\rangle
     =\frac{|\log_{x_0}p|^2}{2R}.
\]
Completing the square, including $p=x_0$ by continuity, we obtain
\begin{equation}\label{eq:offset-sphere}
 \Gamma\subset\exp_{x_0}
 \left\{v\in T_{x_0}M:|v+RN(x_0)|=R\right\}.
\end{equation}
The Euclidean sphere in \Cref{eq:offset-sphere} passes through the
origin of $T_{x_0}M$; its center is $-RN(x_0)$.
\Cref{fig:offset-sphere} depicts this distinction.

\begin{figure}[htbp]
 \centering
 \begin{tikzpicture}[scale=0.95,>=Latex]
  \fill[blue!6] (-2,0) circle (2);
  \draw[thick,blue!65!black] (-2,0) circle (2);
  \draw[->,gray!65] (-4.5,0)--(1.8,0);
  \draw[->,gray!65] (0,-2.35)--(0,2.5);
  \fill (-2,0) circle (1.5pt) node[below left] {$-RN(x_0)$};
  \fill (0,0) circle (1.7pt) node[below right] {$0\in T_{x_0}M$};
  \draw[thick,->] (0,0)--(1.25,0) node[above] {$N(x_0)$};
  \draw[thick] (-2,0)--(-2,2) node[midway,left] {$R$};
  \node at (-2,-1.1) {$U_0$};
  \node[align=center] at (-2,2.65) {A planar section of the tangent-space sphere};
 \end{tikzpicture}
 \caption{In dimension three, the open ball $U_0\subset T_{x_0}M$
 has radius $R$ and center $-RN(x_0)$, so the logarithm origin lies on
 its boundary. At this stage the perimeter support is known to lie on
 $\exp_{x_0}(\partial U_0)$. The phase-recovery argument identifies the occupied
 phase and then proves that the induced metric on the enclosed region
 is Euclidean.}
 \label{fig:offset-sphere}
\end{figure}



\subsection{Recovering the whole ball and its metric}
The three dimension regimes give the same conclusion: for some $o\in M$,
$\Gamma$ is contained in $\exp_o(S)$, where $S$ is a Euclidean sphere
of radius $R$ in $T_oM$. Let $U_0$ be the open ball bounded by $S$ and
put $U=\exp_o(U_0)$. The ball need not be centered at the tangent-space
origin.

The global exponential diffeomorphism maps the two connected components
of $T_oM\setminus S$ onto those of $M\setminus\exp_o(S)$.
Their interior component $U$ is bounded, since its closure is a compact
image. The exterior component is unbounded: vectors escaping to infinity
in the Euclidean exterior have image distances $d(o,\exp_o v)=|v|$
escaping to infinity. Exterior connectedness uses $n\ge2$.
The smooth hypersurface $\exp_o(S)$ has zero ambient volume.

The distributional derivative of $\chi_E$ is supported on $\Gamma$,
so it vanishes on each complementary component. A function with all
weak coordinate derivatives zero is locally constant almost everywhere,
as follows by mollification in coordinate balls. Overlapping balls
then make $\chi_E$ almost everywhere constant on each connected component.
Its two constants belong to $\{0,1\}$. Boundedness of $E$ forces the
exterior constant to be zero, because the open exterior contains
positive-volume balls arbitrarily far from $E$. Since $V>0$ and the
separating sphere has zero volume, the interior constant is one.
Therefore
\begin{equation}\label{eq:phase-recovery}
 E=U\quad\text{up to a set of volume zero}.
\end{equation}
Only containment of the perimeter support in the sphere was needed.

By \Cref{lem:distance-comparison}, the pullback metric
$\widetilde g=(\exp_o)^*g$ dominates the constant Euclidean metric
$g_o$ on all of $T_oM$, hence also on the offset ball $U_0$.
On that ball, \Cref{eq:phase-recovery} gives
\[
 \vol_{\widetilde g}(U_0)=V=\omega_nR^n=\vol_{g_o}(U_0).
\]
The smooth relative density $\sqrt{\det_{g_o}\widetilde g}$ is at least
one. Equality of its integral with $\vol_{g_o}(U_0)$ and continuity
make it identically one on $U_0$. All eigenvalues of $\widetilde g$
relative to $g_o$ are at least one and their product is one.
Each eigenvalue is therefore one, and
\[
 (\exp_o)^*g=g_o\quad\text{on }U_0.
\]
Thus $\exp_o:U_0\to U$ is a Riemannian isometry, proving the necessity
in \Cref{thm:equality}.

For the converse, let $F:B_R(0)\subset\R^n\to U\subset M$ be a
Riemannian isometry onto an open region with its induced metric.
Put $p=F(0)$ and $L=dF_0$. For $|v|<R$, the image of the segment
$s\mapsto sv$ is an ambient geodesic with initial data $(p,Lv)$:
the Levi--Civita connection on an open region is the restriction of
the ambient connection. Geodesic uniqueness gives
\[
 F(v)=\exp_p(Lv).
\]
Consequently $U=B_R(p)$. This formula extends smoothly to the closed
Euclidean ball through the ambient exponential map. Its pullback metric
is Euclidean in the interior and therefore on the boundary by continuity.
The boundary is a smooth sphere of area $s_{n-1}R^{n-1}$, which equals
its ambient perimeter by the divergence formula. Its volume is
$\omega_nR^n$, so equality holds. Perimeter and volume are unchanged
by null modifications. This proves the sufficiency and completes
\Cref{thm:equality}. No flatness outside the equality region is implied.
